\section{Purpose, claim taxonomy, and the fixed-shift target}
\label{sec:scope-target}

\subsection{Four claim tags and three proof levels}

Every major statement in this paper carries one of four logical tags:
\[
\Etag\ \text{established},\qquad
\Htag\ \text{working hypothesis},\qquad
\Ctag\ \text{conditional synthesis},\qquad
\Ntag\ \text{not established}.
\]
The tags answer a different question from the inherited proof levels:
\(\Lzero\) denotes finite algebra, exact model identities, abstract
operator theorems, or countermodels; \(\Lone\) denotes a literal
physical-frame interface preserving the actual arithmetic typing; and
\(\Ltwo\) is reserved for a new growing estimate on the complete
coefficient attached to one prescribed nonzero shift \(h_0\).

\begin{definition}[Promotion rule]
\label{def:promotion}
An averaged-shift, almost-all-scale, frozen-family, random-sign,
support-only, positive-square, or surrogate-carrier estimate is not an
\(\Ltwo\) estimate for the physical fixed-\(h_0\) coefficient unless
an explicit theorem returns it to that coefficient with all losses
and exceptional sectors controlled.
\end{definition}

This paper is an \(\Lone\) synthesis.  A conditional implication may
have an \(\Ltwo\) conclusion while none of its unproved hypotheses is
promoted to an established \(\Ltwo\) theorem.

\subsection{The target correlation}

Fix throughout a nonzero even integer \(h_0\), and let
\(W\in C_c^\infty((0,\infty))\) be fixed.  Write
\[
\mathcal C_{h_0,W}(X)
=
\sum_{n\ge1}\Lambda(n)\Lambda(n+h_0)W(n/X).
\]
If \(\nu_{h_0}(p)\) is the number of residue classes occupied by
\(\{0,-h_0\}\pmod p\), define the usual singular series
\[
\Ssing(h_0)
=
\prod_p
\frac{1-\nu_{h_0}(p)/p}{(1-1/p)^2}.
\]
It is positive for every fixed admissible even \(h_0\).  The centered
defect is the classical prime-pair target
\cite{HardyLittlewood1923,FriedlanderIwaniec2010}:
\begin{equation}
\Delta_{h_0,W}(X)
=
\mathcal C_{h_0,W}(X)
-\Ssing(h_0)X\int_0^\infty W(t)\,dt.
\label{eq:defect}
\end{equation}

\begin{proposition}[\Etag\ Established hard-packet anchor]
\label{prop:anchor}
The TPC fixed-shift local-model reduction gives, for a suitable fixed
\(0<\delta<1/2\),
\begin{equation}
\Delta_{h_0,W}(X)
=
B_{h_0,\delta}(X)
+O_{A,h_0,W,\delta}\!\left(X(\log X)^{-A}\right)
\label{eq:hard-packet-anchor}
\end{equation}
for every fixed \(A>0\), where \(B_{h_0,\delta}\) is an explicitly
defined bilinear hard packet.  Consequently the weighted
Hardy--Littlewood asymptotic is equivalent, within that reduction, to
\[
B_{h_0,\delta}(X)=o(X).
\]
\end{proposition}

This is the entrance to the architecture, not an estimate of the hard
packet.  The purpose of the later TPC layers is to expose its literal
geometry without replacing it by an easier random or averaged object
\cite{WangTPC15,WangTPC18}.

\subsection{Six invariants that no route may drop}
\label{subsec:invariants}

Every hypothesis and every future promotion is audited against:
\begin{enumerate}[label=\textbf{I\arabic*.},leftmargin=2.6em]
\item the literal physical coefficients, including their signs and
  reconstruction multipliers;
\item one prescribed nonzero \(h_0\), not an averaged or subsequently
  chosen shift;
\item physical atomic normalization and exactly one global
  normalization;
\item canonical or losslessly invertible native keys, without
  null-pair inflation;
\item the actual active support, masks, prefixes, polarizations, short
  fibers, and boundary sectors; and
\item a nonduplicated endpoint ledger whose total physical loss is
  strictly below \(1/400\).
\end{enumerate}
Failure of any item prevents a claim of final proof progress.  In
particular, equality at the physical endpoint is a stop.
